%packages permettant d'augmenter le nombre de registres de dimension et donc d'éviter les erreurs de compilation dûs aux packages tikz, pstricks and co
\usepackage{etex}

\usefonttheme{serif}
% option [handout] à ajouter pour bloquer les animations et rendre le doc imprimable
% option [mathserif] permet de retrouver la police habituelle pour les formules mathématiques ; cela ne modifie pas la police du texte

%Thèmes Beamer :
\usetheme{Warsaw}

%\usetheme[left,hideallsubsections]{Paloalto}
%\usetheme{PaloAlto}

%Pour le thème PaloAlto
%\useoutertheme{sidebar}
%\useinnertheme[shadow=true]{rounded}
%\usecolortheme{orchid}
%\usecolortheme{whale}
%\setbeamercolor*{frametitle}{parent=palette primary}
%\setbeamerfont{block title}{size={}}

\usepackage{tipa}%alphabet phonétique internationnal

\usepackage[upright]{fourier}% l'option permet d'avoir les majuscules droites dans les formules mathématiques
\usepackage[utf8x]{inputenc}
%\usepackage[latin1,utf8x]{inputenc}
\usepackage[T1]{fontenc}
\usepackage[english,french]{babel}
\AddThinSpaceBeforeFootnotes % à insérer si on utilise \usepackage[french]{babel}
\FrenchFootnotes % à insérer si on utilise \usepackage[french]{babel}
%s'utilise ainsi : Les Représentants\footnote{ceci est une note} du Peuple Français
\DecimalMathComma %supprime l'espace après la virgule dans un nombre

\usepackage{enumerate} % pour pouvoir changer la numérotation des listes
\usepackage{enumitem}
\frenchbsetup{StandardLists=true}%frenchb ne s'occupera pas des listes
\setenumerate[1]{font=\bfseries\color{blue},label=\arabic*.} % numérotation 1. 2. ...
\setenumerate[2]{font=\bfseries\color{blue},label=\alph*)} % sous-numérotation a) b) ...

\usepackage[scaled=0.875]{helvet}
\renewcommand{\ttdefault}{lmtt}


\usepackage{mathtools,amsmath,amssymb,makeidx}


%redéfinition de fractions, limites, sommes, intégrales, coefficients binomiaux en displaystyle, limites de suites
\let\binomOld\binom
\renewcommand{\binom}{\displaystyle\binomOld}
%\let\fracOld\frac
%\renewcommand{\frac}{\displaystyle\fracOld}
\let\limOld\lim
\renewcommand{\lim}{\displaystyle\limOld}
\newcommand{\limn}{\lim_{n\to +\infty}}
\newcommand{\limm}{\lim_{x\to -\infty}}
\newcommand{\limp}{\lim_{x\to +\infty}}
\newcommand{\limz}{\lim_{x\to 0}}
\newcommand{\limzm}{\lim_{\substack{x \to 0\\ x < 0}}}
\newcommand{\limzp}{\lim_{\substack{x \to 0\\ x > 0}}}
\let\sumOld\sum
\renewcommand{\sum}{\displaystyle\sumOld}
\let\intOld\int
\renewcommand{\int}{\displaystyle\intOld}

%Calligraphie spéciale
\usepackage{mathrsfs}   % Police de maths jolie caligraphie
\newcommand{\calig}[1]{\ensuremath{\mathscr{#1}}}
\newcommand\mtc[1]{\ensuremath{\mathcal{#1}}}


%Nombres complexes
\let\Reold\Re
\renewcommand{\Re}{~\text{Re}~}
\let\Imold\Im
\renewcommand{\Im}{~\text{Im}~}
\newcommand{\ii}{\,\mathrm{i}}

%style des pages
\usepackage{fancybox,cancel}
\usepackage{fancyhdr}
\usepackage{lastpage}
%redéfinition du style plain
\fancypagestyle{plain}{%
\fancyhf{} %vide l'en-tête et le pied de page
%\fancyfoot[C]{\bfseries \thepage / \pageref{LastPage}} %numéro de la page en gras et centré
\fancyfoot[C]{\thepage / \pageref{LastPage}} %numéro de la page centré
\renewcommand{\headrulewidth}{0.2pt}
\renewcommand{\footrulewidth}{0.2pt}}

%___________________________
%===    Raccourcis classe
%------------------------------------------------------
\newcommand\seconde{2\up{nde}\xspace}
\newcommand\premiere{1\up{ère}\xspace}
\newcommand\terminale{T\up{le}\xspace}
\newcommand\stmg{\bsc{Stmg}}
\newcommand\sti{\bsc{Sti2d}}
\newcommand\bat{BAT 1\xspace}
\newcommand\BAT{BAT 2\xspace}
\newcommand\tesspe{TES Spécialité\xspace}

\usepackage{stmaryrd}%\llbracket et \rrbracket % crochets doubles pour intervalles d'entier
%symbole parallèle avec \sslash

\newcommand{\crochets}[2]{\ensuremath{\llbracket #1 ; #2 \rrbracket}}

\newcommand{\intervalleff}[2]{\left[#1\,;#2\right]}
\newcommand{\intervallefo}[2]{\left[#1\,;#2\right[}
\newcommand{\intervalleof}[2]{\left]#1\,;#2\right]}
\newcommand{\intervalleoo}[2]{\left]#1\,;#2\right[}


%Gestion des espaces
%
\newcommand{\pv}{\ensuremath{\: ;}}
\newlength{\EspacePV}
\setlength{\EspacePV}{1em plus 0.5em minus 0.5em}
\newcommand{\qq}{\hspace{\EspacePV} ; \hspace{\EspacePV}}
\newcommand{\qetq}{\hspace{\EspacePV} \text{et} \hspace{\EspacePV}}
\newcommand{\qouq}{\hspace{\EspacePV} \text{ou} \hspace{\EspacePV}}
\newcommand{\qLq}{\hspace{\EspacePV} \Leftarrow \hspace{\EspacePV}}
\newcommand{\qRq}{\hspace{\EspacePV} \Rightarrow \hspace{\EspacePV}}
\newcommand{\qLRq}{\hspace{\EspacePV} \Leftrightarrow \hspace{\EspacePV}}

%TABLEAU
%diminuer la taille des caractères dans un tableau
%Pour aérer un tableau, il est possible de redéfinir l'espacement entre les lignes d'un tableaux et l'espacement entre les colonnes, par exemple :
\usepackage{array}
%\setlength{\tabcolsep}{1cm}
\renewcommand{\arraystretch}{1.5}%augmente la hauteur des lignes des tableaux
%colonnes centrées verticalement et horizontalement permettant d'écrire des paragraphes de largeur fixée du type M{3cm}
\newcolumntype{M}[1]{>{\centering\arraybackslash}m{#1}}%\arraybackslash permet de continuer à utiliser \\ pour le changement de ligne
%couleurs cellules, colonnes, lignes
\usepackage{color,colortbl}
\usepackage{longtable}%permet d'obtenir des tableaux sur plusieurs pages
\usepackage%[table]
{xcolor}
%\usepackage{arydshln}% pour pouvoir ajouter des lignes horizontales en pointillés avec \hdashline au lieu de \hline
%pose des problèmes d'encadré avec alterqcm

%arbres
\usepackage{pstricks,pst-plot,pst-text,pst-tree,pstricks-add}

%tabvar %packages à ajouter aux précédents
\usepackage{pst-eps,pst-fill,pst-node,pst-math}
%\input tabvar %ajouter le fichier tabvar.tex dans le même dossier que le fichier actuel
%à utiliser uniquement lors d'un tableau de signes et de variations


%INTERLIGNES
\usepackage{setspace}
%s'utilise avec \begin{spacing}{''facteur''}
%   […]
%\end{spacing}

%IMAGES
\usepackage{graphicx} %inclure des graphiques
\usepackage{caption}
\usepackage{subfig}

\usepackage{tabularx}
\usepackage{soul} % Pour souligner : \ul
\usepackage{ulem} % Pour souligner double : \uuline
                      % Pour souligner ondulé : \uwave
                      % Pour barrer horizontal : \sout
                      % Pour barrer diagonal : \xout
\usepackage{slashbox}
\usepackage{textcomp}
\usepackage{pst-plot}

\usepackage[np]{numprint}


%ENCADRES
\usepackage{pstricks}
\usepackage{pst-grad}
\usepackage{xkeyval}
\usepackage{pst-coil}
\usepackage{ifthen}
\usepackage{ifpdf}
\usepackage{pst-blur}
\usepackage{bclogo}
%doc ici :
%http://www.tug.org/texlive/Contents/live/texmf-dist/doc/latex/bclogo/bclogo-doc.pdf

%figures tikz
\usepackage{tikz,tikz-3dplot}

\usetikzlibrary{calc,shapes,arrows,plotmarks,lindenmayersystems,decorations,decorations.markings,decorations.pathmorphing,
decorations.pathreplacing,patterns,positioning,decorations.text,shadows,trees}

\usepackage{tkz-base,tkz-fct,tkz-euclide,tkz-tab,tkz-graph}

\usetkzobj{all}


\usepackage{pgfplots}
\pgfplotsset{compat=1.14}


%raccourcis perso
\newcommand\pfr[1]{\fcolorbox{red}{white}{#1}}
\newcommand\coef[1][]{c{\oe}fficient#1\xspace}
\newcommand\abs[1]{\ensuremath{\left\vert #1 \right\vert}}%valeur absolue
\newcommand\Arc[1]{\ensuremath{\wideparen{#1}}}%arc de cercle

%liens hypertexte
%\usepackage[colorlinks=true,linkcolor=black,filecolor=blue,urlcolor=blue]
\usepackage{hyperref}
%\href{mailto:mail@exemple.com}{Mon mail} % Lien email
%\href{http://exemple.com}{Mon site web}  % Lien web
%\href{fichier.pdf}{Mon fichier}          % Lien vers un fichier

\newcommand{\euro}{\texteuro{}~}
\usepackage{eurosym}%avec \EUR{3,12}

%simplification notation norme \norme{}
\newcommand{\norme}[1]{\left\Vert #1\right\Vert}

%Ensembles de nombres
\usepackage{dsfont}
\newcommand{\C}{\mathds C}
\newcommand{\R}{\mathds R}
\newcommand{\Q}{\mathds Q}
\newcommand{\D}{\mathds D}
\newcommand{\Z}{\mathds Z}
\newcommand{\N}{\mathds N}
\newcommand{\U}{\mathds U}
\newcommand{\K}{\mathds K}
\newcommand\Ind{\mathds 1} %= fonction indicatrice
\newcommand\p{\mathds P} %= probabilité
\newcommand\E{\mathds E} % Espérance
\newcommand\V{\mathds V} % Variance
\newcommand{\e}{\mathrm {e}}
\newcommand{\dd}{\,\mathrm{d}}
\newcommand{\pgcd}{\text{pgcd}}
%\newcommand{\si}{\,\text{si}\,}
\newcommand{\sinon}{\,\text{sinon}\,}
\newcommand{\Id}{\text{Id}}
\newcommand{\Vect}{\text{Vect}}

%simplification de la notation de vecteur \vect{}
\newcommand{\vect}[1]{\mathchoice%
{\overrightarrow{\displaystyle\mathstrut#1\,\,}}%
{\overrightarrow{\textstyle\mathstrut#1\,\,}}%
{\overrightarrow{\scriptstyle\mathstrut#1\,\,}}%
{\overrightarrow{\scriptscriptstyle\mathstrut#1\,\,}}}




%Repères
\def\Oij{$\left(\text{O}\pv\vect{\imath},~\vect{\jmath}\right)$\xspace}
\def\Oijk{$\left(\text{O}\pv\vect{\imath},~ \vect{\jmath},~ \vect{k}\right)$\xspace}
\def\Ouv{$\left(\text{O}\pv\vect{u},~\vect{v}\right)$\xspace}
\def\OIJ{$\left(O\pv I\:,\,J\right)$\xspace}





%Symboles Casio dans les programmes
\usepackage{pifont} % ça c'est pour le retour chariot
\newcommand{\RetourChariot}{\Pisymbol{psy}{191}}
\newcommand{\triangleCasio}{
\begin{tikzpicture}[scale=0.2]
\tkzDefPoint(0,0){A}
\tkzDefPoint(1,0){B}
\tkzDefPoint(1,1){C}
\tkzDrawPolygon[fill](A,B,C)
\end{tikzpicture}
}



%Pointillés sur toute la ligne
\usepackage{multido}

\newcommand{\Pointilles}[1][1]{%
\multido{}{#1}{\makebox[\textwidth]{\dotfill}\\[\parskip]
}}
%commandes : \Pointilles ou \Pointilles[4] pour 4 lignes


%textes à trous
\newlength\lgtrou
\newcommand*\trou[1]{%
\settowidth\lgtrou{#1}%
\hspace*{2\lgtrou}}
\setlength\baselineskip{1.2\baselineskip}
%Commande à utiliser : \trou{texte qui sera remplacé par un espace vide}



%%%% Papier millim\'er\'ex*y en cm
\newcommand{\PapierMill}[2]{%
	\begin{pspicture*}(0,0)(#1,#2)
	
	\psgrid[gridwidth=0.3pt,
		gridlabels=0,
		subgriddiv=10,
		subgridcolor=black,
		subgridwidth=0.05pt](0,0)(21,29.7) % il faut \'ere sr
				% que \'e d\'easse, sinon il y a des soucis avec
				% des lignes qui ne vont pas jusqu'au bout
				% \'ecause des arrondis entier de multido
				% pspicture* coupe aux bonnes dimensions
	
	% Grille de 5cm en 5cm
	\psgrid[xunit=5,yunit=5,gridwidth=0.5pt,
		gridlabels=0,
		subgriddiv=0](0,0)(#1,#2)%
		
	\psline[linewidth=1.5pt](0,0)(0,#2)(#1,#2)(#1,0)(0,0)	
	\end{pspicture*}
	}


% Ecrire sur plusieurs colonnes
\usepackage{multicol}
\setlength{\columnseprule}{0.5pt}
%exemple
%\begin{multicols}{3}[Titre sur une seule colonne.]
%   3~colonnes équilibrées, 3~colonnes équilibrées, 3~colonnes équilibrées, 3~colonnes équilibrées
%\end{multicols}
%\begin{multicols}{2}[\section{Titre numéroté.}]
%   blabla sur deux colonnes, c'est plus sérieux. C'est le style qui est généralement utilisé pour écrire des articles.
%saut de colonne forcé : 
%\columnbreak
%djhskjdhjsq
%sdkksqjhd
%\end{multicols}
%Pour ajouter un titre numéroté qui apparaisse sur toute la largeur de la page, il faut utiliser l'option [\section{Titre.}] juste après \begin{multicols}{nb-col}.
%Remarques :
%+ Pour qu'une ligne de séparation apparaisse entre les colonnes, il faut utiliser : \setlength{\columnseprule}{1pt}.

%+ Pour redéfinir la largeur de l'espace inter-colonnes, il faut utiliser \setlength{\columnsep}{30pt}.

%Pour remonter le texte, dans chaque colonne vers le haut : \raggedcolumns qui se tape :\begin{multicols}{2}\raggedcolumns...\columnbreak...\columnbreak\end{multicols}

%Pour supprimer les traits verticaux : \setlength{\columnseprule}{0pt} avant \begin{multicols}{3}...\end{multicols}

%QRcode, codebarre
\usepackage{pst-barcode}
%\begin{pspicture}(2,2)
%	\psbarcode{http://www.latex-howto.be}{eclevel=M}{qrcode}
%\end{pspicture}

%Texte en filigrane
\usepackage{watermark}
%On utilise ensuite les commandes \watermark, \leftwatermark, \rightwatermark ou \thiswatermark qui permettent de définir un filigrane sur toutes les pages, les pages paires, les pages impaires ou juste une page
%Exemple : \thiswatermark {
%\begin{minipage}{0.95\linewidth}
%\vspace{25cm}
%\begin{center}
%\rotatebox{55}{\scalebox{8}{\color[gray]{0.7}\LaTeX}}
%\end{center}
%\end{minipage}
%}

%centrer du texte ou une formule avec moins d'espace autour
\newcommand{\centrer}[1]
{
\smallskip
\centerline{#1}
\smallskip
}


%Permettre la même numérotation beamer
\newcounter{MonPetitCompteur}
\setbeamertemplate{section in toc}{%
\edef\MaDefTemp{\noexpand\setcounter{MonPetitCompteur}{\inserttocsectionnumber}}\MaDefTemp%
\Roman{MonPetitCompteur} - \inserttocsection\par}

\setbeamertemplate{subsection in toc}{%
\leavevmode\leftskip=1.5em%
\edef\MaDefTemp{\noexpand\setcounter{MonPetitCompteur}{\inserttocsubsectionnumber}}\MaDefTemp%
\arabic{MonPetitCompteur}) \inserttocsubsection\par}

%\usepackage{mathtools} %pour l'alignement des termes d'une matrice

% Mise en forme des algorithmes
\usepackage[french,boxed,titlenumbered,lined,longend]{algorithm2e}
  \SetKwIF {Si}{SinonSi}{Sinon}{si}{alors}{sinon\_si}{alors}{fin~si}
 \SetKwFor{Tq}{tant\_que~}{~faire~}{fin~tant\_que}
 \SetKwFor{PourCh}{pour\_chaque }{ faire }{fin pour\_chaque}
 \SetKwInput{Sortie}{Sortie}
  \SetKwInput{Entree}{Entrée}
\newcommand{\Algocmd}[1]{\textsf{\textsc{\textbf{#1}}}}\SetKwSty{Algocmd}
  \newcommand{\AlgCommentaire}[1]{\textsl{\small  #1}} 
  

\usepackage{pas-algo} %package de Stéphane Pasquet
\usepackage{tcolorbox}  
  
%%%%%%%%%%%%%%%%%%%%%%%%%%%%%%%%%%
%Pages de départ de la présentation
%%%%%%%%%%%%%%%%%%%%%%%%%%%%%%%%%%

\AtBeginSection{
   \begin{frame}%{\textcolor{red}{\shadowbox{Plan}}}
   \Large%
   \tableofcontents[currentsection,currentsubsection]
   \end{frame}}

\AtBeginSubsection{  
  \begin{frame}%{\textcolor{red}{\shadowbox{Plan}}}
   \Large%
 \tableofcontents[currentsection,currentsubsection]
   \end{frame}}




%%%%%%%%%%%%%%%%%%%%%%%%%%%%%%%%%%%%%%%%%%%%%%%%%%%%%%%%%%%%%%%%%%%%%%%%%%%%%%%

\pagecolor{white}
%%%%%%%%%%%%%%%%%%%%%%%%%%%%%%%%%%%%%%%%%%%%%%%%%%%%%%%%%%%%%%%%%%%%%%%%%%%%%%%
%%%%%%%%%%%%%%%%%%%%%%%%%%%%%%%%%%%%%%%%%%%%%%%%%%%%%%%%%%%%%%%%%%%%%%%%%%%%%%%

%___________________________
%===    TABLEUR
%------------------------------------------------------
\usepackage{pas-tableur}%package de Stéphane Pasquet
\usepackage{xstring}
\usepackage{xkeyval}



%%%%%%%%%Autres commandes
%symbole pour variable aléatoire qui suit une loi
\newcommand{\suit}{\hookrightarrow}



%%%%%%%%%%%%%%%%%%%%%%%%%%%%%%%%%%%%%%%%%%%%%%%%%%%%%%%%%%%%%%%%%%%%%%%%%%%
%%%%%%%%%%%%%%%%%%%%%%%%%%%%%%%%%%%%%%%%%%%%%%%%%%%%%%%%%%%%%%%%%%%%%%%%%%%
%%%%%%%%%%%%%%%%%%%%%%%%%%% Les environnements du cours %%%%%%%%%%%%%%%%%%%
%%%%%%%%%%%%%%%%%%%%%%%%%%%%%%%%%%%%%%%%%%%%%%%%%%%%%%%%%%%%%%%%%%%%%%%%%%%
%%%%%%%%%%%%%%%%%%%%%%%%%%%%%%%%%%%%%%%%%%%%%%%%%%%%%%%%%%%%%%%%%%%%%%%%%%%
%definition des couleurs
\definecolor{fondpaille}{cmyk}{0,0,0.1,0}%\pagecolor{fondpaille}
\definecolor{gris}{rgb}{0.7,0.7,0.7}
\definecolor{rouge}{rgb}{1,0,0}
\definecolor{bleu}{rgb}{0,0,1}
\definecolor{vert}{rgb}{0,1,0}
\definecolor{deficolor}{HTML}{2D9AFF}
\definecolor{backdeficolor}{HTML}{EDEDED}%{036DD0}%dégradé bleu{666666}%dégradé gris
\definecolor{theocolor}{HTML}{C10CC7}%{HTML}{036DD0}%F4404D%rouge
\definecolor{backtheocolor}{HTML}{D3D3D3}
\definecolor{methcolor}{HTML}{008800}%12BB05}
\definecolor{backmethcolor}{HTML}{FFFACD}
\definecolor{backilluscolor}{HTML}{EDEDED}
\definecolor{sectioncolor}{HTML}{C10CC7}%{B2B2B2}%vert : {HTML}{008800}%{HTML}{2D9AFF}
\definecolor{subsectioncolor}{HTML}{C10CC7}%{B2B2B2}%vert : {HTML}{008800}%{rgb}{0.5,0,0}
\definecolor{subsubsectioncolor}{HTML}{C10CC7}
\definecolor{engcolor}{HTML}{D4D7FE}
\definecolor{exocolor}{rgb}{0,0.6,0}
\definecolor{exosoltitlecolor}{rgb}{0,0.6,0}
\definecolor{titlecolor}{rgb}{1,1,1}

%commande pour enlever les couleurs avant impression
\newcommand{\nocolor}
{\pagecolor{white}
\definecolor{gris}{rgb}{0.7,0.7,0.7}
\definecolor{rouge}{rgb}{0,0,0}
\definecolor{bleu}{rgb}{0,0,0}
\definecolor{vert}{rgb}{0,0,0}
\definecolor{deficolor}{HTML}{B2B2B2}
\definecolor{backdeficolor}{HTML}{EEEEEE}%{036DD0}%dégradé bleu{666666}%dégradé gris
\definecolor{theocolor}{HTML}{B2B2B2}
\definecolor{backtheocolor}{HTML}{EEEEEE}
\definecolor{methcolor}{HTML}{B2B2B2}
\definecolor{backmethcolor}{HTML}{EEEEEE}
\definecolor{backilluscolor}{HTML}{EEEEEE}
\definecolor{sectioncolor}{HTML}{B2B2B2}
\definecolor{subsectioncolor}{HTML}{B2B2B2}
\definecolor{subsubsectioncolor}{HTML}{B2B2B2}
\definecolor{engcolor}{HTML}{EEEEEE}
\definecolor{exocolor}{HTML}{3B3838}
\definecolor{exosoltitlecolor}{rgb}{0,0,0}
\definecolor{titlecolor}{rgb}{0,0,0}
}


%%%%%%%%%%%%%%%%%%%%%%%%%%%%%%%%%%%%%%%%%%%%%%%%%%%%%%%%%%%%%%%%%%%%%%%%%%%%%%%
%Encadrés pour Propriétés, Théorème, Définitions, exemples, exercices

\usepackage{environ}%pour pouvoir utiliser la commande \NewEnviron

%___________________________
%===    Propriété avec ou sans s et avec ou sans titre
%------------------------------------------------------
%
\newcounter{prop}
\NewEnviron{Prop}[2][]{
\smallskip
\refstepcounter{prop}
\begin{tikzpicture}[node distance=0 cm]
\node[fill=theocolor,%rounded corners=5pt,
anchor=south west] (theorem) at (0,0)
{\textcolor{titlecolor}{\bfseries Propriété#1~\theprop~:~#2}};
\node[%draw,
drop shadow,%color=deficolor,very thick,
fill=backtheocolor,%rounded corners=5pt,
anchor=north west] at(0,-0.02)
{\black%\parbox
\begin{minipage}{\columnwidth-12pt}%{
\BODY%}
\end{minipage}};
\end{tikzpicture}
\smallskip
}

%___________________________
%===    Notation avec ou sans s et avec ou sans titre
%------------------------------------------------------
%
\newcounter{notation}
\NewEnviron{Notation}[2][]{
\smallskip
\refstepcounter{notation}
\begin{tikzpicture}[node distance=0 cm]
\node[fill=red,%rounded corners=5pt,
anchor=south west] (theorem) at (0,0)
{\textcolor{titlecolor}{\bfseries Notation#1~\thenotation~:~#2}};
\node[%draw,
drop shadow,%color=deficolor,very thick,
fill=red!50,%rounded corners=5pt,
anchor=north west] at(0,-0.02)
{\black%\parbox
\begin{minipage}{\columnwidth-12pt}%{
\BODY%}
\end{minipage}};
\end{tikzpicture}
\smallskip
}


%___________________________
%===    Conséquence avec ou sans s et avec ou sans titre
%------------------------------------------------------
%
\newcounter{cons}
\NewEnviron{Cons}[2][]{
\smallskip
\refstepcounter{cons}
\begin{tikzpicture}[node distance=0 cm]
\node[fill=theocolor,%rounded corners=5pt,
anchor=south west] (theorem) at (0,0)
{\textcolor{titlecolor}{\bfseries Conséquence#1~\thecons~:~#2}};
\node[%draw,
drop shadow,%color=deficolor,very thick,
fill=backtheocolor,%rounded corners=5pt,
anchor=north west] at(0,-0.02)
{\black%\parbox
\begin{minipage}{\columnwidth-12pt}%{
\BODY%}
\end{minipage}};
\end{tikzpicture}
\smallskip
}


%___________________________
%===    Corolaire avec ou sans titre
%------------------------------------------------------
%
\newcounter{cor}
\NewEnviron{Cor}[1][]{
\smallskip
\refstepcounter{cor}
\begin{tikzpicture}[node distance=0 cm]
\node[fill=theocolor,%rounded corners=5pt,
anchor=south west] (theorem) at (0,0)
{\textcolor{titlecolor}{\bfseries Corollaire~\thecor~:~#1}};
\node[%draw,
drop shadow,%color=deficolor,very thick,
fill=backtheocolor,%rounded corners=5pt,
anchor=north west] at(0,-0.02)
{\black%\parbox
\begin{minipage}{\columnwidth-12pt}%{
\BODY%}
\end{minipage}};
\end{tikzpicture}
\smallskip
}


%___________________________
%===    Théorème avec ou sans titre
%------------------------------------------------------
%
\newcounter{thm}
\NewEnviron{Thm}[1][]{
\smallskip
\refstepcounter{thm}
\begin{tikzpicture}[node distance=0 cm]
\node[fill=theocolor,%rounded corners=5pt,
anchor=south west] (theorem) at (0,0)
{\textcolor{titlecolor}{\bfseries Théorème~\thethm~:~#1}};
\node[%draw,
drop shadow,%color=deficolor,very thick,
fill=backtheocolor,%rounded corners=5pt,
anchor=north west] at(0,-0.02)
{\black%\parbox
\begin{minipage}{\columnwidth-12pt}%{
\BODY%}
\end{minipage}};
\end{tikzpicture}
\smallskip
}


%___________________________
%===    Cadre arrondi coloré
%------------------------------------------------------
%
\NewEnviron{CadreColor}{
\medskip
\begin{tikzpicture}[node distance=0 cm]
\node[draw,drop shadow,color=theocolor,very thick,fill=backtheocolor,rounded corners=5pt,anchor=north west] at(0,-0.02)
{\black\parbox{\linewidth-12pt}{\BODY}};
\end{tikzpicture}
%\medskip
}


%___________________________
%===    Cadre arrondi blanc
%------------------------------------------------------
%
\NewEnviron{Cadre}{
\medskip
\begin{tikzpicture}[node distance=0 cm]
\node[draw,very thick,rounded corners=5pt,anchor=north west] at(0,-0.02)
{\black\parbox{\linewidth-12pt}{\BODY}};
\end{tikzpicture}
%\medskip
}


%_______________________________________________
%===    Cadre présentation algo et modules
%------------------------------------------------------
%

%utilise le package tcolorbox
\newtcolorbox{mybox}[2][]{%enhanced,
colback=red!5!white,
colframe=red!75!black,fonttitle=\bfseries,
colbacktitle=red!85!black,
%attach boxed title to top left={xshift=2mm,yshift=-2mm},
title=#2,#1}

\newtcolorbox{module}[2][]{%enhanced,
colback=red!5!white,
colframe=red!75!black,fonttitle=\bfseries,
colbacktitle=red!85!black,
%attach boxed title to bottom right={xshift=-5mm,yshift=+2mm},
title=#2,#1}

\NewEnviron{CadreAlgo}[2]{
\begin{minipage}{#1} %largeur de la minipage
\begin{mybox}[colback=blue!10]{#2} %Titre du cadre
\BODY
\end{mybox}
\end{minipage}
}

\NewEnviron{CadreModule}[1]{
\begin{module}[colback=blue!10]{#1}%Catégorie du module
\BODY %Titre de module
\end{module}
}


%___________________________
%===    Règle(s) avec ou sans s et avec sans titre
%------------------------------------------------------
%
\newcounter{regle}
\NewEnviron{Regle}[2][]{
\smallskip
\refstepcounter{regle}
\begin{tikzpicture}[node distance=0 cm]
\node[fill=theocolor,%rounded corners=5pt,
anchor=south west] (theorem) at (0,0)
{\textcolor{titlecolor}{\bfseries Règle#1~\theregle~:~#2}};
\node[%draw,
drop shadow,%color=deficolor,very thick,
fill=backdeficolor,%rounded corners=5pt,
anchor=north west] at(0,-0.02)
{\black%\parbox
\begin{minipage}{\columnwidth-12pt}%{
\BODY%}
\end{minipage}};
\end{tikzpicture}
\smallskip
}

%___________________________
%===    Vocabulaire avec ou sans titre
%------------------------------------------------------
%
\newcounter{vocab}
\NewEnviron{Vocab}[1][]{
\smallskip
\refstepcounter{vocab}
\begin{tikzpicture}[node distance=0 cm]
\node[fill=theocolor,%rounded corners=5pt,
anchor=south west] (theorem) at (0,0)
{\textcolor{titlecolor}{\bfseries Vocabulaire~\thevocab~:~#1}};
\node[%draw,
drop shadow,%color=deficolor,very thick,
fill=backdeficolor,%rounded corners=5pt,
anchor=north west] at(0,-0.02)
{\black%\parbox
\begin{minipage}{\columnwidth-12pt}%{
\BODY%}
\end{minipage}};
\end{tikzpicture}
\smallskip
}


%___________________________
%===    Définition avec ou sans s et avec sans titre
%------------------------------------------------------
%
\newcounter{defi}
\NewEnviron{Defi}[2][]{
\smallskip
\refstepcounter{defi}
\begin{tikzpicture}[node distance=0 cm]
\node[fill=theocolor,%rounded corners=5pt,
anchor=south west] (theorem) at (0,0)
{\textcolor{titlecolor}{\bfseries Définition#1~\thedefi~:~#2}};
\node[%draw,
drop shadow,%color=deficolor,very thick,
fill=backdeficolor,%rounded corners=5pt,
anchor=north west] at(0,-0.02)
{\black%\parbox
\begin{minipage}{\columnwidth-12pt}%{
\BODY%}
\end{minipage}};
\end{tikzpicture}
\smallskip
}

%___________________________
%===    Méthode avec ou sans s et avec sans titre
%------------------------------------------------------
%
\newcounter{methode}
\NewEnviron{Methode}[2][]{
\smallskip
\refstepcounter{methode}
\begin{tikzpicture}[node distance=0 cm]
\node[fill=theocolor,%rounded corners=5pt,
anchor=south west] (theorem) at (0,0)
{\textcolor{titlecolor}{\bfseries Méthode#1~\themethode~:~#2}};
\node[%draw,
drop shadow,%color=deficolor,very thick,
fill=backmethcolor,%rounded corners=5pt,
anchor=north west] at(0,-0.02)
{\black%\parbox
\begin{minipage}{\columnwidth-12pt}%{
\BODY%}
\end{minipage}};
\end{tikzpicture}
\smallskip
}



%___________________________
%===    Exemple avec ou sans s et avec ou sans titre
%------------------------------------------------------
%
\newcounter{exemple}
\NewEnviron{Exemple}[2][]{
\smallskip
\refstepcounter{exemple}
\textbf{\large{Exemple#1~\theexemple~:~#2}}\par
\BODY
\smallskip
}

%___________________________
%===    Remarque avec ou sans s
%------------------------------------------------------
%
\newcounter{rmq}
\NewEnviron{Rmq}[1][]{
\smallskip
\refstepcounter{rmq}
\textbf{\large{Remarque#1~\thermq~:}}\par
\BODY
\smallskip
}

%___________________________
%===    Remarques numérotées R1, R2, etc...
%------------------------------------------------------
%
\newcounter{rem}\newcommand{\rem}{\refstepcounter{rem}\textbf{R \therem \ :}\xspace}

%___________________________
%===    Exercices du contrôle numérotés
%------------------------------------------------------
%
\newcounter{exercice}
\NewEnviron{Exercice}[1][]{
\smallskip
\refstepcounter{exercice}\textbf{\large{Exercice \theexercice \ :}}\hfill \textbf{#1}\par
\BODY
\smallskip
}


%___________________________
%===    Environnements CNED
%------------------------------------------------------
%
\newcounter{propcned}
\NewEnviron{PropCNED}[2][]{\refstepcounter{propcned}
\smallskip
\begin{tikzpicture}[node distance=0 cm]
\node[fill=theocolor,rounded corners=5pt,anchor=south west] (theorem) at (0,0)
{\textcolor{titlecolor}{\bfseries Propriété#1 \thepropcned ~:~#2}};
\node[draw,drop shadow,color=theocolor,very thick,fill=backtheocolor,rounded corners=5pt,anchor=north west] at(0,-0.02)
{\black\parbox{\columnwidth-12pt}{\BODY}};
\end{tikzpicture}
\smallskip
}

\newcounter{defcned}
\NewEnviron{DefCNED}[2][]{\refstepcounter{defcned}
\smallskip
\begin{tikzpicture}[node distance=0 cm]
\node[fill=theocolor,rounded corners=5pt,anchor=south west] (theorem) at (0,0)
{\textcolor{titlecolor}{\bfseries Définition#1 \thedefcned ~:~#2}};
\node[draw,drop shadow,color=theocolor,very thick,fill=backtheocolor,rounded corners=5pt,anchor=north west] at(0,-0.02)
{\black\parbox{\columnwidth-12pt}{\BODY}};
\end{tikzpicture}
\smallskip
}

\newcounter{thmcned}
\NewEnviron{ThmCNED}[2][]{\refstepcounter{thmcned}
\smallskip
\begin{tikzpicture}[node distance=0 cm]
\node[fill=theocolor,rounded corners=5pt,anchor=south west] (theorem) at (0,0)
{\textcolor{titlecolor}{\bfseries Théorème#1 \thethmcned ~:~#2}};
\node[draw,drop shadow,color=theocolor,very thick,fill=backtheocolor,rounded corners=5pt,anchor=north west] at(0,-0.02)
{\black\parbox{\columnwidth-12pt}{\BODY}};
\end{tikzpicture}
\smallskip
}

\newcounter{exemplecned}
\NewEnviron{ExempleCNED}[2][]{\refstepcounter{exemplecned}
\smallskip
\textbf{\large{Exemple#1 \theexemplecned~:~#2}}\par
\BODY
\smallskip
}

\newcounter{rmqcned}
\NewEnviron{RmqCNED}[2][]{\refstepcounter{rmqcned}
\smallskip
\textbf{\large{Remarque#1 \thermqcned~:~#2}}\par
\BODY
\smallskip
}

%___________________________
%===    Exercices du contrôle numérotés en Anglais
%------------------------------------------------------
%
\newcounter{exercise}
\NewEnviron{Exercise}[1][]{
\smallskip
\refstepcounter{exercise}\textbf{\large{Exercise \#\theexercise:}}\hfill \textbf{#1}\par
\BODY
\smallskip
}

%___________________________
%===    Exercices non numérotés
%------------------------------------------------------
%
\NewEnviron{Exo}[1][]{
\smallskip
\textbf{\large{Exercice #1 \ :}}\par
\BODY
\smallskip
}

%___________________________
%===    Exercices non numérotés en Anglais
%------------------------------------------------------
%
\NewEnviron{Exoa}[1][]{
\smallskip
\textbf{\large{Exercise #1 \ :}}\par
\BODY
\smallskip
}

%___________________________
%===    Démonstration
%------------------------------------------------------
\newcounter{demo}
\NewEnviron{Demo}[1][]{%
\smallskip
\refstepcounter{demo}
\textit{\textbf{Démonstration \thedemo #1}}\par
\BODY
\strut\hfill$\square$
\smallskip
}

%___________________________
%===    Commandes perso
%------------------------------------------------------




%Commande chapitre
\newcommand{\chapitre}[1]{
\begin{tcolorbox}[colframe=sectioncolor,colback=sectioncolor]
\begin{center}
\vspace*{0.5em}
\textcolor{titlecolor}{\Large \textsc{\textbf{#1}}}
\vspace*{0.5em}
\end{center}
\end{tcolorbox}
\bigskip
}

%Pour les fiches : commande de Cécile
\newcommand{\Fiche}[2]{%
\begin{tcolorbox}[colframe=sectioncolor,colback=sectioncolor]
\begin{center}
\vspace*{0.5em}
\textcolor{blue}{\Large \textsc{\textbf{Fiche~#1\ :\ #2}}}
\vspace*{0.5em}
\end{center}
\end{tcolorbox}
}%

\pagecolor{white}%couleur du fond de page


%Pour pouvoir utiliser l'environnement verbatim
\usepackage{verbatim}

%Panneau danger (nécessite le package pstricks)
\def\danger{\begingroup
\psset{unit=1ex}%
\begin{pspicture}(0,0)(3,3)
\pspolygon[linearc=0.2,linewidth=0.12,linecolor=red](0,0)(1.5,2.6)(3,0)
\psellipse*(1.5,1.33)(0.14,0.75)\pscircle*(1.5,0.3){0.15}\end{pspicture}
\endgroup}%

\newcommand{\cad}{c'est-à-dire }


%Fonctions hyperboliques

\newcommand{\ch}{\text{ch}}
\newcommand{\sh}{\text{sh}}
\renewcommand{\th}{\text{th}}
\newcommand{\argch}{\text{argch}}
\newcommand{\argsh}{\text{argsh}}
\newcommand{\argth}{\text{argth}}


%___________________________
%===    Nouvelles commandes pour documents venant de Sesamath
%------------------------------------------------------
%
\definecolor{CyanTikz40}{cmyk}{.4,0,0,0}
\definecolor{CyanTikz20}{cmyk}{.2,0,0,0}

\definecolor{B1prime}                {cmyk}{0.00, 1.00, 0.00, 0.50}
\definecolor{H1prime}                {cmyk}{0.50, 0.00, 1.00, 0.00}

\definecolor{FootFonctionColor}{cmyk}{0.50, 0.00, 0.00, 0.00}
\definecolor{FootGeometrieColor}{cmyk}{0.40, 0.40, 0.00, 0.00}
\definecolor{FootStatistiqueColor}{cmyk}{0.30, 0.48, 0.00, 0.10}
\definecolor{FootStatistiqueOLDColor}{cmyk}{0.48, 0.30, 0.10, 0.00}
\definecolor{FootStatistique*Color}{cmyk}{0.20, 0.00, 0.00, 0.00}
\definecolor{ActiviteFootColor}{cmyk}{0.50, 0.00, 0.25, 0.00}
\definecolor{CoursFootColor}{cmyk}{0.15, 0.00, 0.00, 0.03}
\definecolor{ExoBaseFootColor}{cmyk}{0.00, 0.25, 0.50, 0.00}
\definecolor{ExoApprFootColor}{cmyk}{0.00, 0.25, 0.50, 0.00}
%\colorlet{ConnFootColor}{F2}
\definecolor{TPFootColor}{cmyk}{0.00, 0.30, 0.00, 0.10}
\definecolor{RecreationFootColor}{cmyk}{0.20, 0.00, 0.50, 0.05}

\definecolor{Blanc}             {cmyk}{0.00, 0.00, 0.00, 0.00}
\definecolor{Gris1}             {cmyk}{0.00, 0.00, 0.00, 0.20}
\definecolor{Gris2}             {cmyk}{0.00, 0.00, 0.00, 0.40}
\definecolor{Gris3}             {cmyk}{0.00, 0.00, 0.00, 0.50}
\definecolor{Noir}              {cmyk}{0.00, 0.00, 0.00, 1.00}
\definecolor{A1}              {cmyk}{0.33, 1.00, 0.00, 0.40}
\definecolor{F1}              {cmyk}{0.00, 1.00, 1.00, 0.00}
\definecolor{C1}              {cmyk}{0.00, 1.00, 0.00, 0.50}
\definecolor{G1}              {cmyk}{0.00, 0.00, 0.00, 0.20}
\definecolor{D1}              {cmyk}{0.00, 0.22, 0.49, 0.69}%bitume
\definecolor{J1}              {cmyk}{0.00, 0.34, 1.00, 0.02}%orangé


\tikzstyle{general}=[line width=0.3mm, x=1cm, y=1cm,line cap=round, line join=round]
\tikzstyle{quadrillage}=[line width=0.3mm, color=CyanTikz40]
\tikzstyle{quadrillageNIV2}=[line width=0.3mm, color=CyanTikz20]
\tikzstyle{quadrillage55}=[line width=0.3mm, color=CyanTikz40, xstep=0.5, ystep=0.5]
\tikzstyle{cote}=[line width=0.3mm, <->,>=stealth]
\tikzstyle{epais}=[line width=0.5mm, line cap=butt]
\tikzstyle{tres epais}=[line width=0.8mm, line cap=butt]
\tikzstyle{axe}=[line width=0.6mm, ->,>=stealth, color=blue, line cap=rect]
\newcommand{\quadrillageSeyes}[2]{\draw[line width=0.3mm, color=A1!10, ystep=0.2, xstep=0.8] #1 grid #2;
\draw[line width=0.3mm, color=A1!30, xstep=0.8, ystep=0.8] #1 grid #2; }
\newcommand{\axeX}[4][0]{\draw[axe] (#2,#1)--(#3,#1); \foreach \x in {#4} {\draw (\x,#1) node {\small $+$}; \draw (\x,#1) node[below] {\small $\x$};}}
\newcommand{\axeY}[4][0]{\draw[axe] (#1,#2)--(#1,#3); \foreach \y in {#4} {\draw (#1, \y) node {\small $+$}; \draw (#1, \y) node[left] {\small $\y$};}}
\newcommand{\axeOI}[3][0]{\draw[axe] (#2,#1)--(#3,#1);  \draw (1,#1) node {\small $+$}; \draw (1,#1) node[below] {\small $I$};}
\newcommand{\axeOJ}[3][0]{\draw[axe] (#1,#2)--(#1,#3); \draw (#1, 1) node {\small $+$}; \draw (#1, 1) node[left] {\small $J$};}
\newcommand{\axeXgraduation}[2][0]{\foreach \x in {#2} {\draw (\x,#1) node {\small $+$};}}
\newcommand{\axeYgraduation}[2][0]{\foreach \y in {#2} {\draw (#1, \y) node {\small $+$}; }}
\newcommand{\origine}{\draw (0,0) node[below left] {\small $0$};}
\newcommand{\origineO}{\draw (0,0) node[below left] {$O$};}
\newcommand{\point}[4]{\draw (#1,#2) node[#4] {$#3$};}
\newcommand{\pointGraphique}[4]{\draw (#1,#2) node[#4] {$#3$};
\draw (#1,#2) node[scale=1.5] {$\times$};}
\newcommand{\pointFigure}[4]{\draw (#1,#2) node[#4] {$#3$};
\draw (#1,#2) node[scale=1.5] {$\times$};}
\newcommand{\pointC}[3]{\draw (#1) node[#3] {$#2$};}
\newcommand{\pointCGraphique}[3]{\draw (#1) node[#3] {$#2$};
\draw (#1) node[scale=1.5] {$\times$};}
\newcommand{\pointCFigure}[3]{\draw (#1) node[#3] {$#2$};
\draw (#1) node[scale=1.5] {$\times$};}


\newcommand{\expco}[1]{\e^{#1}}
\newcommand{\expcopi}[2]{\e^{\ii \frac{{#1\pi}}{#2}}}
\newcommand{\expcopim}[2]{\e^{-\ii \frac{{#1\pi}}{#2}}}
\newcommand{\expcopimul}[3]{\e^{{ #1\ii} \frac{{#2\pi}}{#3}}}
\newcommand{\cj}{\mathrm{j}}
\newcommand{\ci}{\ii}


%augmenter l'espace au-dessus ou en-dessous d'une fraction
\makeatletter
\newcommand*\Strut[1][1]{%
  \leavevmode
  \vrule \@height #1\ht\strutbox
         \@depth #1\dp\strutbox
         \@width\z@
}
\newcommand*\TopStrut[1][1]{%
  \leavevmode
  \vrule \@height #1\ht\strutbox
         \@depth \z@
         \@width \z@
}
\newcommand*\BotStrut[1][1]{%
  \leavevmode
  \vrule \@height \z@
         \@depth #1\dp\strutbox
         \@width \z@
}
\makeatother

%exemple
%Résoudre: $  \dfrac{ \TopStrut 3 x -1} { \BotStrut x + 2} < 3$.



%%%%%%%%%%%%%%%%%%%%%%%%%%%%%%%%
%Symbole pour les probabilités
\DeclareMathOperator{\Prob}{\mathbb{P}}
%%%%%%%%%%%%%%%%%%%%%%%%%%%%%%%%%%%%%


%%%%%%%%%%%%%%%%%%%%%%%%%%%%%%%%%%%%%
%\usepackage{bbding}
%permet d'utiliser \Checkmark et \XSolidBrush

%%%%%%%%%%%%%%%%%%%%%%%%%%%%%%%%%%%%%

%%%%%%%%%%%%%%%%%%%%%%%%%%%%%%%%%%%%%%%%%%%%%%%%%%%%%%%%%%%%%%
%%%%%%%%%%%%%%%%%%%%%%%%%%%% Scratch %%%%%%%%%%%%%%%%%%%%%%%%%%
%%%%%%%%%%%%%%%%%%%%%%%%%%%%%%%%%%%%%%%%%%%%%%%%%%%%%%%%%%%%%%

\usepackage{scratch3}


%%%%%%%%%%%%%%%%%%%%%%%%%%%%%%%%%%%%%%%%%%%%%%%%%%%%%%%%%%%%%%
%%%%%%%%%%%%%%%%%%%%%%%%%%%% Python %%%%%%%%%%%%%%%%%%%%%%%%%%
%%%%%%%%%%%%%%%%%%%%%%%%%%%%%%%%%%%%%%%%%%%%%%%%%%%%%%%%%%%%%%

%nécessite l'image correspondante
\newcommand{\python}{\raisebox{-0.8em}{\includegraphics[scale=0.3]{python-logo}}}

%Nécessaire pour l'environnement lstlisting
\usepackage{listings}

%\begin{lstlisting}[language=Python]
%# Calcul de la factorielle
%def factorielle(x):
%	if x < 2:
%		return 1
%	else:
%		return x * factorielle(x-1)
%str(5) + "! = " + str(factorielle(5))
%\end{lstlisting}


\definecolor{purple2}{RGB}{153,0,153} % there's actually no standard purple
\definecolor{green2}{RGB}{0,153,0} % a darker green
\lstset{%
language=Python, % the language
basicstyle=\normalsize\ttfamily, % size of the fonts for the code
% Color settings to match IDLE style
keywordstyle=\color{red}, % core keywords
keywordstyle={[2]\color{purple2}}, % built-ins
stringstyle=\color{green2},
commentstyle=\color{blue},
upquote=true, % requires textcomp
literate=
  {á}{{\'a}}1 {é}{{\'e}}1 {í}{{\'i}}1 {ó}{{\'o}}1 {ú}{{\'u}}1
  {Á}{{\'A}}1 {É}{{\'E}}1 {Í}{{\'I}}1 {Ó}{{\'O}}1 {Ú}{{\'U}}1
  {à}{{\`a}}1 {è}{{\`e}}1 {ì}{{\`i}}1 {ò}{{\`o}}1 {ù}{{\`u}}1
  {À}{{\`A}}1 {È}{{\'E}}1 {Ì}{{\`I}}1 {Ò}{{\`O}}1 {Ù}{{\`U}}1
  {ä}{{\"a}}1 {ë}{{\"e}}1 {ï}{{\"i}}1 {ö}{{\"o}}1 {ü}{{\"u}}1
  {Ä}{{\"A}}1 {Ë}{{\"E}}1 {Ï}{{\"I}}1 {Ö}{{\"O}}1 {Ü}{{\"U}}1
  {â}{{\^a}}1 {ê}{{\^e}}1 {î}{{\^i}}1 {ô}{{\^o}}1 {û}{{\^u}}1
  {Â}{{\^A}}1 {Ê}{{\^E}}1 {Î}{{\^I}}1 {Ô}{{\^O}}1 {Û}{{\^U}}1
  {œ}{{\oe}}1 {Œ}{{\OE}}1 {æ}{{\ae}}1 {Æ}{{\AE}}1 {ß}{{\ss}}1
  {ű}{{\H{u}}}1 {Ű}{{\H{U}}}1 {ő}{{\H{o}}}1 {Ő}{{\H{O}}}1
  {ç}{{\c c}}1 {Ç}{{\c C}}1 {ø}{{\o}}1 {å}{{\r a}}1 {Å}{{\r A}}1
  {€}{{\EUR}}1 {£}{{\pounds}}1,
frame=single,% style du cadre - none pour ne pas encadrer
frameround=fttt,
numbers=none,%left,% numéros des lignes à gauche - none pour ne pas numéroter
numberstyle=\small,% on peut mettre \tiny ou \normalsize (taille par défaut)
firstnumber=0,% numéro de la première ligne
tabsize=6,% largeur des tabulations, 8 par défaut
showtabs=false,% permet de voir les tabulations quand true
xleftmargin=5pt,% marge gauche
xrightmargin=5pt,% marge droite qui peut rester à 0
framexleftmargin=0pt,% pour inclure les numéros dans le cadre : mettre 20pt
showspaces=false,% n'affiche pas de caractères pour remplacer les espaces
showstringspaces=false,% n'affiche pas de caractères pour remplacer les espaces dans les string
breaklines=true,%passage à ligne automatique
}


 %%%%%%%%%%%%%%%%%%%%%%%%%%%%%%%%%%%%%%%%%%%%%%%%%%%%%%%%%%%%%%%%%%%%%%%%%%%%%%%% 
%%% ~ Arduino Language - Arduino IDE Colors ~                                  %%%
%%%                                                                            %%%
%%% Kyle Rocha-Brownell | 10/2/2017 | No Licence                               %%%
%%% -------------------------------------------------------------------------- %%%
%%%                                                                            %%%
%%% Place this file in your working directory (next to the latex file you're   %%%
%%% working on).  To add it to your project, place:                            %%%
%%%     %%%%%%%%%%%%%%%%%%%%%%%%%%%%%%%%%%%%%%%%%%%%%%%%%%%%%%%%%%%%%%%%%%%%%%%%%%%%%%%% 
%%% ~ Arduino Language - Arduino IDE Colors ~                                  %%%
%%%                                                                            %%%
%%% Kyle Rocha-Brownell | 10/2/2017 | No Licence                               %%%
%%% -------------------------------------------------------------------------- %%%
%%%                                                                            %%%
%%% Place this file in your working directory (next to the latex file you're   %%%
%%% working on).  To add it to your project, place:                            %%%
%%%     %%%%%%%%%%%%%%%%%%%%%%%%%%%%%%%%%%%%%%%%%%%%%%%%%%%%%%%%%%%%%%%%%%%%%%%%%%%%%%%% 
%%% ~ Arduino Language - Arduino IDE Colors ~                                  %%%
%%%                                                                            %%%
%%% Kyle Rocha-Brownell | 10/2/2017 | No Licence                               %%%
%%% -------------------------------------------------------------------------- %%%
%%%                                                                            %%%
%%% Place this file in your working directory (next to the latex file you're   %%%
%%% working on).  To add it to your project, place:                            %%%
%%%     %%%%%%%%%%%%%%%%%%%%%%%%%%%%%%%%%%%%%%%%%%%%%%%%%%%%%%%%%%%%%%%%%%%%%%%%%%%%%%%% 
%%% ~ Arduino Language - Arduino IDE Colors ~                                  %%%
%%%                                                                            %%%
%%% Kyle Rocha-Brownell | 10/2/2017 | No Licence                               %%%
%%% -------------------------------------------------------------------------- %%%
%%%                                                                            %%%
%%% Place this file in your working directory (next to the latex file you're   %%%
%%% working on).  To add it to your project, place:                            %%%
%%%    \input{arduinoLanguage.tex}                                             %%%
%%% somewhere before \begin{document} in your latex file.                      %%%
%%%                                                                            %%%
%%% In your document, place your arduino code between:                         %%%
%%%   \begin{lstlisting}[language=Arduino]                                     %%%
%%% and:                                                                       %%%
%%%   \end{lstlisting}                                                         %%%
%%%                                                                            %%%
%%% Or create your own style to add non-built-in functions and variables.      %%%
%%%                                                                            %%%
 %%%%%%%%%%%%%%%%%%%%%%%%%%%%%%%%%%%%%%%%%%%%%%%%%%%%%%%%%%%%%%%%%%%%%%%%%%%%%%%% 

\usepackage{color}
\usepackage{listings}    
\usepackage{courier}

%%% Define Custom IDE Colors %%%
\definecolor{arduinoGreen}    {rgb} {0.17, 0.43, 0.01}
\definecolor{arduinoGrey}     {rgb} {0.47, 0.47, 0.33}
\definecolor{arduinoOrange}   {rgb} {0.8 , 0.4 , 0   }
\definecolor{arduinoBlue}     {rgb} {0.01, 0.61, 0.98}
\definecolor{arduinoDarkBlue} {rgb} {0.0 , 0.2 , 0.5 }

%%% Define Arduino Language %%%
\lstdefinelanguage{Arduino}{
  language=C++, % begin with default C++ settings 
%
%
  %%% Keyword Color Group 1 %%%  (called KEYWORD3 by arduino)
  keywordstyle=\color{arduinoGreen},   
  deletekeywords={  % remove all arduino keywords that might be in c++
                break, case, override, final, continue, default, do, else, for, 
                if, return, goto, switch, throw, try, while, setup, loop, export, 
                not, or, and, xor, include, define, elif, else, error, if, ifdef, 
                ifndef, pragma, warning,
                HIGH, LOW, INPUT, INPUT_PULLUP, OUTPUT, DEC, BIN, HEX, OCT, PI, 
                HALF_PI, TWO_PI, LSBFIRST, MSBFIRST, CHANGE, FALLING, RISING, 
                DEFAULT, EXTERNAL, INTERNAL, INTERNAL1V1, INTERNAL2V56, LED_BUILTIN, 
                LED_BUILTIN_RX, LED_BUILTIN_TX, DIGITAL_MESSAGE, FIRMATA_STRING, 
                ANALOG_MESSAGE, REPORT_DIGITAL, REPORT_ANALOG, SET_PIN_MODE, 
                SYSTEM_RESET, SYSEX_START, auto, int8_t, int16_t, int32_t, int64_t, 
                uint8_t, uint16_t, uint32_t, uint64_t, char16_t, char32_t, operator, 
                enum, delete, bool, boolean, byte, char, const, false, float, double, 
                null, NULL, int, long, new, private, protected, public, short, 
                signed, static, volatile, String, void, true, unsigned, word, array, 
                sizeof, dynamic_cast, typedef, const_cast, struct, static_cast, union, 
                friend, extern, class, reinterpret_cast, register, explicit, inline, 
                _Bool, complex, _Complex, _Imaginary, atomic_bool, atomic_char, 
                atomic_schar, atomic_uchar, atomic_short, atomic_ushort, atomic_int, 
                atomic_uint, atomic_long, atomic_ulong, atomic_llong, atomic_ullong, 
                virtual, PROGMEM,
                Serial, Serial1, Serial2, Serial3, SerialUSB, Keyboard, Mouse,
                abs, acos, asin, atan, atan2, ceil, constrain, cos, degrees, exp, 
                floor, log, map, max, min, radians, random, randomSeed, round, sin, 
                sq, sqrt, tan, pow, bitRead, bitWrite, bitSet, bitClear, bit, 
                highByte, lowByte, analogReference, analogRead, 
                analogReadResolution, analogWrite, analogWriteResolution, 
                attachInterrupt, detachInterrupt, digitalPinToInterrupt, delay, 
                delayMicroseconds, digitalWrite, digitalRead, interrupts, millis, 
                micros, noInterrupts, noTone, pinMode, pulseIn, pulseInLong, shiftIn, 
                shiftOut, tone, yield, Stream, begin, end, peek, read, print, 
                println, available, availableForWrite, flush, setTimeout, find, 
                findUntil, parseInt, parseFloat, readBytes, readBytesUntil, readString, 
                readStringUntil, trim, toUpperCase, toLowerCase, charAt, compareTo, 
                concat, endsWith, startsWith, equals, equalsIgnoreCase, getBytes, 
                indexOf, lastIndexOf, length, replace, setCharAt, substring, 
                toCharArray, toInt, press, release, releaseAll, accept, click, move, 
                isPressed, isAlphaNumeric, isAlpha, isAscii, isWhitespace, isControl, 
                isDigit, isGraph, isLowerCase, isPrintable, isPunct, isSpace, 
                isUpperCase, isHexadecimalDigit, 
                }, 
  morekeywords={   % add arduino structures to group 1
                break, case, override, final, continue, default, do, else, for, 
                if, return, goto, switch, throw, try, while, setup, loop, export, 
                not, or, and, xor, include, define, elif, else, error, if, ifdef, 
                ifndef, pragma, warning,
                }, 
% 
%
  %%% Keyword Color Group 2 %%%  (called LITERAL1 by arduino)
  keywordstyle=[2]\color{arduinoBlue},   
  keywords=[2]{   % add variables and dataTypes as 2nd group  
                HIGH, LOW, INPUT, INPUT_PULLUP, OUTPUT, DEC, BIN, HEX, OCT, PI, 
                HALF_PI, TWO_PI, LSBFIRST, MSBFIRST, CHANGE, FALLING, RISING, 
                DEFAULT, EXTERNAL, INTERNAL, INTERNAL1V1, INTERNAL2V56, LED_BUILTIN, 
                LED_BUILTIN_RX, LED_BUILTIN_TX, DIGITAL_MESSAGE, FIRMATA_STRING, 
                ANALOG_MESSAGE, REPORT_DIGITAL, REPORT_ANALOG, SET_PIN_MODE, 
                SYSTEM_RESET, SYSEX_START, auto, int8_t, int16_t, int32_t, int64_t, 
                uint8_t, uint16_t, uint32_t, uint64_t, char16_t, char32_t, operator, 
                enum, delete, bool, boolean, byte, char, const, false, float, double, 
                null, NULL, int, long, new, private, protected, public, short, 
                signed, static, volatile, String, void, true, unsigned, word, array, 
                sizeof, dynamic_cast, typedef, const_cast, struct, static_cast, union, 
                friend, extern, class, reinterpret_cast, register, explicit, inline, 
                _Bool, complex, _Complex, _Imaginary, atomic_bool, atomic_char, 
                atomic_schar, atomic_uchar, atomic_short, atomic_ushort, atomic_int, 
                atomic_uint, atomic_long, atomic_ulong, atomic_llong, atomic_ullong, 
                virtual, PROGMEM,
                },  
% 
%
  %%% Keyword Color Group 3 %%%  (called KEYWORD1 by arduino)
  keywordstyle=[3]\bfseries\color{arduinoOrange},
  keywords=[3]{  % add built-in functions as a 3rd group
                Serial, Serial1, Serial2, Serial3, SerialUSB, Keyboard, Mouse,
                },      
%
%
  %%% Keyword Color Group 4 %%%  (called KEYWORD2 by arduino)
  keywordstyle=[4]\color{arduinoOrange},
  keywords=[4]{  % add more built-in functions as a 4th group
                abs, acos, asin, atan, atan2, ceil, constrain, cos, degrees, exp, 
                floor, log, map, max, min, radians, random, randomSeed, round, sin, 
                sq, sqrt, tan, pow, bitRead, bitWrite, bitSet, bitClear, bit, 
                highByte, lowByte, analogReference, analogRead, 
                analogReadResolution, analogWrite, analogWriteResolution, 
                attachInterrupt, detachInterrupt, digitalPinToInterrupt, delay, 
                delayMicroseconds, digitalWrite, digitalRead, interrupts, millis, 
                micros, noInterrupts, noTone, pinMode, pulseIn, pulseInLong, shiftIn, 
                shiftOut, tone, yield, Stream, begin, end, peek, read, print, 
                println, available, availableForWrite, flush, setTimeout, find, 
                findUntil, parseInt, parseFloat, readBytes, readBytesUntil, readString, 
                readStringUntil, trim, toUpperCase, toLowerCase, charAt, compareTo, 
                concat, endsWith, startsWith, equals, equalsIgnoreCase, getBytes, 
                indexOf, lastIndexOf, length, replace, setCharAt, substring, 
                toCharArray, toInt, press, release, releaseAll, accept, click, move, 
                isPressed, isAlphaNumeric, isAlpha, isAscii, isWhitespace, isControl, 
                isDigit, isGraph, isLowerCase, isPrintable, isPunct, isSpace, 
                isUpperCase, isHexadecimalDigit, 
                },      
%
%
  %%% Set Other Colors %%%
  stringstyle=\color{arduinoDarkBlue},    
  commentstyle=\color{arduinoGrey},    
%          
%   
  %%%% Line Numbering %%%%
   numbers=left,                    
  numbersep=5pt,                   
  numberstyle=\color{arduinoGrey},    
  %stepnumber=2,                      % show every 2 line numbers
%
%
  %%%% Code Box Style %%%%
  breaklines=true,                    % wordwrapping
  tabsize=2,         
  basicstyle=\ttfamily  
}                                             %%%
%%% somewhere before \begin{document} in your latex file.                      %%%
%%%                                                                            %%%
%%% In your document, place your arduino code between:                         %%%
%%%   \begin{lstlisting}[language=Arduino]                                     %%%
%%% and:                                                                       %%%
%%%   \end{lstlisting}                                                         %%%
%%%                                                                            %%%
%%% Or create your own style to add non-built-in functions and variables.      %%%
%%%                                                                            %%%
 %%%%%%%%%%%%%%%%%%%%%%%%%%%%%%%%%%%%%%%%%%%%%%%%%%%%%%%%%%%%%%%%%%%%%%%%%%%%%%%% 

\usepackage{color}
\usepackage{listings}    
\usepackage{courier}

%%% Define Custom IDE Colors %%%
\definecolor{arduinoGreen}    {rgb} {0.17, 0.43, 0.01}
\definecolor{arduinoGrey}     {rgb} {0.47, 0.47, 0.33}
\definecolor{arduinoOrange}   {rgb} {0.8 , 0.4 , 0   }
\definecolor{arduinoBlue}     {rgb} {0.01, 0.61, 0.98}
\definecolor{arduinoDarkBlue} {rgb} {0.0 , 0.2 , 0.5 }

%%% Define Arduino Language %%%
\lstdefinelanguage{Arduino}{
  language=C++, % begin with default C++ settings 
%
%
  %%% Keyword Color Group 1 %%%  (called KEYWORD3 by arduino)
  keywordstyle=\color{arduinoGreen},   
  deletekeywords={  % remove all arduino keywords that might be in c++
                break, case, override, final, continue, default, do, else, for, 
                if, return, goto, switch, throw, try, while, setup, loop, export, 
                not, or, and, xor, include, define, elif, else, error, if, ifdef, 
                ifndef, pragma, warning,
                HIGH, LOW, INPUT, INPUT_PULLUP, OUTPUT, DEC, BIN, HEX, OCT, PI, 
                HALF_PI, TWO_PI, LSBFIRST, MSBFIRST, CHANGE, FALLING, RISING, 
                DEFAULT, EXTERNAL, INTERNAL, INTERNAL1V1, INTERNAL2V56, LED_BUILTIN, 
                LED_BUILTIN_RX, LED_BUILTIN_TX, DIGITAL_MESSAGE, FIRMATA_STRING, 
                ANALOG_MESSAGE, REPORT_DIGITAL, REPORT_ANALOG, SET_PIN_MODE, 
                SYSTEM_RESET, SYSEX_START, auto, int8_t, int16_t, int32_t, int64_t, 
                uint8_t, uint16_t, uint32_t, uint64_t, char16_t, char32_t, operator, 
                enum, delete, bool, boolean, byte, char, const, false, float, double, 
                null, NULL, int, long, new, private, protected, public, short, 
                signed, static, volatile, String, void, true, unsigned, word, array, 
                sizeof, dynamic_cast, typedef, const_cast, struct, static_cast, union, 
                friend, extern, class, reinterpret_cast, register, explicit, inline, 
                _Bool, complex, _Complex, _Imaginary, atomic_bool, atomic_char, 
                atomic_schar, atomic_uchar, atomic_short, atomic_ushort, atomic_int, 
                atomic_uint, atomic_long, atomic_ulong, atomic_llong, atomic_ullong, 
                virtual, PROGMEM,
                Serial, Serial1, Serial2, Serial3, SerialUSB, Keyboard, Mouse,
                abs, acos, asin, atan, atan2, ceil, constrain, cos, degrees, exp, 
                floor, log, map, max, min, radians, random, randomSeed, round, sin, 
                sq, sqrt, tan, pow, bitRead, bitWrite, bitSet, bitClear, bit, 
                highByte, lowByte, analogReference, analogRead, 
                analogReadResolution, analogWrite, analogWriteResolution, 
                attachInterrupt, detachInterrupt, digitalPinToInterrupt, delay, 
                delayMicroseconds, digitalWrite, digitalRead, interrupts, millis, 
                micros, noInterrupts, noTone, pinMode, pulseIn, pulseInLong, shiftIn, 
                shiftOut, tone, yield, Stream, begin, end, peek, read, print, 
                println, available, availableForWrite, flush, setTimeout, find, 
                findUntil, parseInt, parseFloat, readBytes, readBytesUntil, readString, 
                readStringUntil, trim, toUpperCase, toLowerCase, charAt, compareTo, 
                concat, endsWith, startsWith, equals, equalsIgnoreCase, getBytes, 
                indexOf, lastIndexOf, length, replace, setCharAt, substring, 
                toCharArray, toInt, press, release, releaseAll, accept, click, move, 
                isPressed, isAlphaNumeric, isAlpha, isAscii, isWhitespace, isControl, 
                isDigit, isGraph, isLowerCase, isPrintable, isPunct, isSpace, 
                isUpperCase, isHexadecimalDigit, 
                }, 
  morekeywords={   % add arduino structures to group 1
                break, case, override, final, continue, default, do, else, for, 
                if, return, goto, switch, throw, try, while, setup, loop, export, 
                not, or, and, xor, include, define, elif, else, error, if, ifdef, 
                ifndef, pragma, warning,
                }, 
% 
%
  %%% Keyword Color Group 2 %%%  (called LITERAL1 by arduino)
  keywordstyle=[2]\color{arduinoBlue},   
  keywords=[2]{   % add variables and dataTypes as 2nd group  
                HIGH, LOW, INPUT, INPUT_PULLUP, OUTPUT, DEC, BIN, HEX, OCT, PI, 
                HALF_PI, TWO_PI, LSBFIRST, MSBFIRST, CHANGE, FALLING, RISING, 
                DEFAULT, EXTERNAL, INTERNAL, INTERNAL1V1, INTERNAL2V56, LED_BUILTIN, 
                LED_BUILTIN_RX, LED_BUILTIN_TX, DIGITAL_MESSAGE, FIRMATA_STRING, 
                ANALOG_MESSAGE, REPORT_DIGITAL, REPORT_ANALOG, SET_PIN_MODE, 
                SYSTEM_RESET, SYSEX_START, auto, int8_t, int16_t, int32_t, int64_t, 
                uint8_t, uint16_t, uint32_t, uint64_t, char16_t, char32_t, operator, 
                enum, delete, bool, boolean, byte, char, const, false, float, double, 
                null, NULL, int, long, new, private, protected, public, short, 
                signed, static, volatile, String, void, true, unsigned, word, array, 
                sizeof, dynamic_cast, typedef, const_cast, struct, static_cast, union, 
                friend, extern, class, reinterpret_cast, register, explicit, inline, 
                _Bool, complex, _Complex, _Imaginary, atomic_bool, atomic_char, 
                atomic_schar, atomic_uchar, atomic_short, atomic_ushort, atomic_int, 
                atomic_uint, atomic_long, atomic_ulong, atomic_llong, atomic_ullong, 
                virtual, PROGMEM,
                },  
% 
%
  %%% Keyword Color Group 3 %%%  (called KEYWORD1 by arduino)
  keywordstyle=[3]\bfseries\color{arduinoOrange},
  keywords=[3]{  % add built-in functions as a 3rd group
                Serial, Serial1, Serial2, Serial3, SerialUSB, Keyboard, Mouse,
                },      
%
%
  %%% Keyword Color Group 4 %%%  (called KEYWORD2 by arduino)
  keywordstyle=[4]\color{arduinoOrange},
  keywords=[4]{  % add more built-in functions as a 4th group
                abs, acos, asin, atan, atan2, ceil, constrain, cos, degrees, exp, 
                floor, log, map, max, min, radians, random, randomSeed, round, sin, 
                sq, sqrt, tan, pow, bitRead, bitWrite, bitSet, bitClear, bit, 
                highByte, lowByte, analogReference, analogRead, 
                analogReadResolution, analogWrite, analogWriteResolution, 
                attachInterrupt, detachInterrupt, digitalPinToInterrupt, delay, 
                delayMicroseconds, digitalWrite, digitalRead, interrupts, millis, 
                micros, noInterrupts, noTone, pinMode, pulseIn, pulseInLong, shiftIn, 
                shiftOut, tone, yield, Stream, begin, end, peek, read, print, 
                println, available, availableForWrite, flush, setTimeout, find, 
                findUntil, parseInt, parseFloat, readBytes, readBytesUntil, readString, 
                readStringUntil, trim, toUpperCase, toLowerCase, charAt, compareTo, 
                concat, endsWith, startsWith, equals, equalsIgnoreCase, getBytes, 
                indexOf, lastIndexOf, length, replace, setCharAt, substring, 
                toCharArray, toInt, press, release, releaseAll, accept, click, move, 
                isPressed, isAlphaNumeric, isAlpha, isAscii, isWhitespace, isControl, 
                isDigit, isGraph, isLowerCase, isPrintable, isPunct, isSpace, 
                isUpperCase, isHexadecimalDigit, 
                },      
%
%
  %%% Set Other Colors %%%
  stringstyle=\color{arduinoDarkBlue},    
  commentstyle=\color{arduinoGrey},    
%          
%   
  %%%% Line Numbering %%%%
   numbers=left,                    
  numbersep=5pt,                   
  numberstyle=\color{arduinoGrey},    
  %stepnumber=2,                      % show every 2 line numbers
%
%
  %%%% Code Box Style %%%%
  breaklines=true,                    % wordwrapping
  tabsize=2,         
  basicstyle=\ttfamily  
}                                             %%%
%%% somewhere before \begin{document} in your latex file.                      %%%
%%%                                                                            %%%
%%% In your document, place your arduino code between:                         %%%
%%%   \begin{lstlisting}[language=Arduino]                                     %%%
%%% and:                                                                       %%%
%%%   \end{lstlisting}                                                         %%%
%%%                                                                            %%%
%%% Or create your own style to add non-built-in functions and variables.      %%%
%%%                                                                            %%%
 %%%%%%%%%%%%%%%%%%%%%%%%%%%%%%%%%%%%%%%%%%%%%%%%%%%%%%%%%%%%%%%%%%%%%%%%%%%%%%%% 

\usepackage{color}
\usepackage{listings}    
\usepackage{courier}

%%% Define Custom IDE Colors %%%
\definecolor{arduinoGreen}    {rgb} {0.17, 0.43, 0.01}
\definecolor{arduinoGrey}     {rgb} {0.47, 0.47, 0.33}
\definecolor{arduinoOrange}   {rgb} {0.8 , 0.4 , 0   }
\definecolor{arduinoBlue}     {rgb} {0.01, 0.61, 0.98}
\definecolor{arduinoDarkBlue} {rgb} {0.0 , 0.2 , 0.5 }

%%% Define Arduino Language %%%
\lstdefinelanguage{Arduino}{
  language=C++, % begin with default C++ settings 
%
%
  %%% Keyword Color Group 1 %%%  (called KEYWORD3 by arduino)
  keywordstyle=\color{arduinoGreen},   
  deletekeywords={  % remove all arduino keywords that might be in c++
                break, case, override, final, continue, default, do, else, for, 
                if, return, goto, switch, throw, try, while, setup, loop, export, 
                not, or, and, xor, include, define, elif, else, error, if, ifdef, 
                ifndef, pragma, warning,
                HIGH, LOW, INPUT, INPUT_PULLUP, OUTPUT, DEC, BIN, HEX, OCT, PI, 
                HALF_PI, TWO_PI, LSBFIRST, MSBFIRST, CHANGE, FALLING, RISING, 
                DEFAULT, EXTERNAL, INTERNAL, INTERNAL1V1, INTERNAL2V56, LED_BUILTIN, 
                LED_BUILTIN_RX, LED_BUILTIN_TX, DIGITAL_MESSAGE, FIRMATA_STRING, 
                ANALOG_MESSAGE, REPORT_DIGITAL, REPORT_ANALOG, SET_PIN_MODE, 
                SYSTEM_RESET, SYSEX_START, auto, int8_t, int16_t, int32_t, int64_t, 
                uint8_t, uint16_t, uint32_t, uint64_t, char16_t, char32_t, operator, 
                enum, delete, bool, boolean, byte, char, const, false, float, double, 
                null, NULL, int, long, new, private, protected, public, short, 
                signed, static, volatile, String, void, true, unsigned, word, array, 
                sizeof, dynamic_cast, typedef, const_cast, struct, static_cast, union, 
                friend, extern, class, reinterpret_cast, register, explicit, inline, 
                _Bool, complex, _Complex, _Imaginary, atomic_bool, atomic_char, 
                atomic_schar, atomic_uchar, atomic_short, atomic_ushort, atomic_int, 
                atomic_uint, atomic_long, atomic_ulong, atomic_llong, atomic_ullong, 
                virtual, PROGMEM,
                Serial, Serial1, Serial2, Serial3, SerialUSB, Keyboard, Mouse,
                abs, acos, asin, atan, atan2, ceil, constrain, cos, degrees, exp, 
                floor, log, map, max, min, radians, random, randomSeed, round, sin, 
                sq, sqrt, tan, pow, bitRead, bitWrite, bitSet, bitClear, bit, 
                highByte, lowByte, analogReference, analogRead, 
                analogReadResolution, analogWrite, analogWriteResolution, 
                attachInterrupt, detachInterrupt, digitalPinToInterrupt, delay, 
                delayMicroseconds, digitalWrite, digitalRead, interrupts, millis, 
                micros, noInterrupts, noTone, pinMode, pulseIn, pulseInLong, shiftIn, 
                shiftOut, tone, yield, Stream, begin, end, peek, read, print, 
                println, available, availableForWrite, flush, setTimeout, find, 
                findUntil, parseInt, parseFloat, readBytes, readBytesUntil, readString, 
                readStringUntil, trim, toUpperCase, toLowerCase, charAt, compareTo, 
                concat, endsWith, startsWith, equals, equalsIgnoreCase, getBytes, 
                indexOf, lastIndexOf, length, replace, setCharAt, substring, 
                toCharArray, toInt, press, release, releaseAll, accept, click, move, 
                isPressed, isAlphaNumeric, isAlpha, isAscii, isWhitespace, isControl, 
                isDigit, isGraph, isLowerCase, isPrintable, isPunct, isSpace, 
                isUpperCase, isHexadecimalDigit, 
                }, 
  morekeywords={   % add arduino structures to group 1
                break, case, override, final, continue, default, do, else, for, 
                if, return, goto, switch, throw, try, while, setup, loop, export, 
                not, or, and, xor, include, define, elif, else, error, if, ifdef, 
                ifndef, pragma, warning,
                }, 
% 
%
  %%% Keyword Color Group 2 %%%  (called LITERAL1 by arduino)
  keywordstyle=[2]\color{arduinoBlue},   
  keywords=[2]{   % add variables and dataTypes as 2nd group  
                HIGH, LOW, INPUT, INPUT_PULLUP, OUTPUT, DEC, BIN, HEX, OCT, PI, 
                HALF_PI, TWO_PI, LSBFIRST, MSBFIRST, CHANGE, FALLING, RISING, 
                DEFAULT, EXTERNAL, INTERNAL, INTERNAL1V1, INTERNAL2V56, LED_BUILTIN, 
                LED_BUILTIN_RX, LED_BUILTIN_TX, DIGITAL_MESSAGE, FIRMATA_STRING, 
                ANALOG_MESSAGE, REPORT_DIGITAL, REPORT_ANALOG, SET_PIN_MODE, 
                SYSTEM_RESET, SYSEX_START, auto, int8_t, int16_t, int32_t, int64_t, 
                uint8_t, uint16_t, uint32_t, uint64_t, char16_t, char32_t, operator, 
                enum, delete, bool, boolean, byte, char, const, false, float, double, 
                null, NULL, int, long, new, private, protected, public, short, 
                signed, static, volatile, String, void, true, unsigned, word, array, 
                sizeof, dynamic_cast, typedef, const_cast, struct, static_cast, union, 
                friend, extern, class, reinterpret_cast, register, explicit, inline, 
                _Bool, complex, _Complex, _Imaginary, atomic_bool, atomic_char, 
                atomic_schar, atomic_uchar, atomic_short, atomic_ushort, atomic_int, 
                atomic_uint, atomic_long, atomic_ulong, atomic_llong, atomic_ullong, 
                virtual, PROGMEM,
                },  
% 
%
  %%% Keyword Color Group 3 %%%  (called KEYWORD1 by arduino)
  keywordstyle=[3]\bfseries\color{arduinoOrange},
  keywords=[3]{  % add built-in functions as a 3rd group
                Serial, Serial1, Serial2, Serial3, SerialUSB, Keyboard, Mouse,
                },      
%
%
  %%% Keyword Color Group 4 %%%  (called KEYWORD2 by arduino)
  keywordstyle=[4]\color{arduinoOrange},
  keywords=[4]{  % add more built-in functions as a 4th group
                abs, acos, asin, atan, atan2, ceil, constrain, cos, degrees, exp, 
                floor, log, map, max, min, radians, random, randomSeed, round, sin, 
                sq, sqrt, tan, pow, bitRead, bitWrite, bitSet, bitClear, bit, 
                highByte, lowByte, analogReference, analogRead, 
                analogReadResolution, analogWrite, analogWriteResolution, 
                attachInterrupt, detachInterrupt, digitalPinToInterrupt, delay, 
                delayMicroseconds, digitalWrite, digitalRead, interrupts, millis, 
                micros, noInterrupts, noTone, pinMode, pulseIn, pulseInLong, shiftIn, 
                shiftOut, tone, yield, Stream, begin, end, peek, read, print, 
                println, available, availableForWrite, flush, setTimeout, find, 
                findUntil, parseInt, parseFloat, readBytes, readBytesUntil, readString, 
                readStringUntil, trim, toUpperCase, toLowerCase, charAt, compareTo, 
                concat, endsWith, startsWith, equals, equalsIgnoreCase, getBytes, 
                indexOf, lastIndexOf, length, replace, setCharAt, substring, 
                toCharArray, toInt, press, release, releaseAll, accept, click, move, 
                isPressed, isAlphaNumeric, isAlpha, isAscii, isWhitespace, isControl, 
                isDigit, isGraph, isLowerCase, isPrintable, isPunct, isSpace, 
                isUpperCase, isHexadecimalDigit, 
                },      
%
%
  %%% Set Other Colors %%%
  stringstyle=\color{arduinoDarkBlue},    
  commentstyle=\color{arduinoGrey},    
%          
%   
  %%%% Line Numbering %%%%
   numbers=left,                    
  numbersep=5pt,                   
  numberstyle=\color{arduinoGrey},    
  %stepnumber=2,                      % show every 2 line numbers
%
%
  %%%% Code Box Style %%%%
  breaklines=true,                    % wordwrapping
  tabsize=2,         
  basicstyle=\ttfamily  
}                                             %%%
%%% somewhere before \begin{document} in your latex file.                      %%%
%%%                                                                            %%%
%%% In your document, place your arduino code between:                         %%%
%%%   \begin{lstlisting}[language=Arduino]                                     %%%
%%% and:                                                                       %%%
%%%   \end{lstlisting}                                                         %%%
%%%                                                                            %%%
%%% Or create your own style to add non-built-in functions and variables.      %%%
%%%                                                                            %%%
 %%%%%%%%%%%%%%%%%%%%%%%%%%%%%%%%%%%%%%%%%%%%%%%%%%%%%%%%%%%%%%%%%%%%%%%%%%%%%%%% 

\usepackage{color}
\usepackage{courier}

%%% Define Custom IDE Colors %%%
\definecolor{arduinoGreen}    {rgb} {0.17, 0.43, 0.01}
\definecolor{arduinoGrey}     {rgb} {0.47, 0.47, 0.33}
\definecolor{arduinoOrange}   {rgb} {0.8 , 0.4 , 0   }
\definecolor{arduinoBlue}     {rgb} {0.01, 0.61, 0.98}
\definecolor{arduinoDarkBlue} {rgb} {0.0 , 0.2 , 0.5 }

%%% Define Arduino Language %%%
\lstdefinelanguage{Arduino}{
  language=C++, % begin with default C++ settings 
%
%
  %%% Keyword Color Group 1 %%%  (called KEYWORD3 by arduino)
  keywordstyle=\bfseries\color{arduinoGreen},   
  deletekeywords={  % remove all arduino keywords that might be in c++
                break, case, override, final, continue, default, do, else, for, 
                if, return, goto, switch, throw, try, while, setup, loop, export, 
                not, or, and, xor, include, define, elif, else, error, if, ifdef, 
                ifndef, pragma, warning,
                HIGH, LOW, INPUT, INPUT_PULLUP, OUTPUT, DEC, BIN, HEX, OCT, PI, 
                HALF_PI, TWO_PI, LSBFIRST, MSBFIRST, CHANGE, FALLING, RISING, 
                DEFAULT, EXTERNAL, INTERNAL, INTERNAL1V1, INTERNAL2V56, LED_BUILTIN, 
                LED_BUILTIN_RX, LED_BUILTIN_TX, DIGITAL_MESSAGE, FIRMATA_STRING, 
                ANALOG_MESSAGE, REPORT_DIGITAL, REPORT_ANALOG, SET_PIN_MODE, 
                SYSTEM_RESET, SYSEX_START, auto, int8_t, int16_t, int32_t, int64_t, 
                uint8_t, uint16_t, uint32_t, uint64_t, char16_t, char32_t, operator, 
                enum, delete, bool, boolean, byte, char, const, false, float, double, 
                null, NULL, int, long, new, private, protected, public, short, 
                signed, static, volatile, String, void, true, unsigned, word, array, 
                sizeof, dynamic_cast, typedef, const_cast, struct, static_cast, union, 
                friend, extern, class, reinterpret_cast, register, explicit, inline, 
                _Bool, complex, _Complex, _Imaginary, atomic_bool, atomic_char, 
                atomic_schar, atomic_uchar, atomic_short, atomic_ushort, atomic_int, 
                atomic_uint, atomic_long, atomic_ulong, atomic_llong, atomic_ullong, 
                virtual, PROGMEM,
                Serial, Serial1, Serial2, Serial3, SerialUSB, Keyboard, Mouse,
                abs, acos, asin, atan, atan2, ceil, constrain, cos, degrees, exp, 
                floor, log, map, max, min, radians, random, randomSeed, round, sin, 
                sq, sqrt, tan, pow, bitRead, bitWrite, bitSet, bitClear, bit, 
                highByte, lowByte, analogReference, analogRead, 
                analogReadResolution, analogWrite, analogWriteResolution, 
                attachInterrupt, detachInterrupt, digitalPinToInterrupt, delay, 
                delayMicroseconds, digitalWrite, digitalRead, interrupts, millis, 
                micros, noInterrupts, noTone, pinMode, pulseIn, pulseInLong, shiftIn, 
                shiftOut, tone, yield, Stream, begin, end, peek, read, print, 
                println, available, availableForWrite, flush, setTimeout, find, 
                findUntil, parseInt, parseFloat, readBytes, readBytesUntil, readString, 
                readStringUntil, trim, toUpperCase, toLowerCase, charAt, compareTo, 
                concat, endsWith, startsWith, equals, equalsIgnoreCase, getBytes, 
                indexOf, lastIndexOf, length, replace, setCharAt, substring, 
                toCharArray, toInt, press, release, releaseAll, accept, click, move, 
                isPressed, isAlphaNumeric, isAlpha, isAscii, isWhitespace, isControl, 
                isDigit, isGraph, isLowerCase, isPrintable, isPunct, isSpace, 
                isUpperCase, isHexadecimalDigit, 
                }, 
  morekeywords={   % add arduino structures to group 1
                break, case, override, final, continue, default, do, else, for, 
                if, return, goto, switch, throw, try, while, setup, loop, export, 
                not, or, and, xor, include, define, elif, else, error, if, ifdef, 
                ifndef, pragma, warning,
                }, 
% 
%
  %%% Keyword Color Group 2 %%%  (called LITERAL1 by arduino)
  keywordstyle=[2]\bfseries\color{arduinoBlue},   
  keywords=[2]{   % add variables and dataTypes as 2nd group  
                HIGH, LOW, INPUT, INPUT_PULLUP, OUTPUT, DEC, BIN, HEX, OCT, PI, 
                HALF_PI, TWO_PI, LSBFIRST, MSBFIRST, CHANGE, FALLING, RISING, 
                DEFAULT, EXTERNAL, INTERNAL, INTERNAL1V1, INTERNAL2V56, LED_BUILTIN, 
                LED_BUILTIN_RX, LED_BUILTIN_TX, DIGITAL_MESSAGE, FIRMATA_STRING, 
                ANALOG_MESSAGE, REPORT_DIGITAL, REPORT_ANALOG, SET_PIN_MODE, 
                SYSTEM_RESET, SYSEX_START, auto, int8_t, int16_t, int32_t, int64_t, 
                uint8_t, uint16_t, uint32_t, uint64_t, char16_t, char32_t, operator, 
                enum, delete, bool, boolean, byte, char, const, false, float, double, 
                null, NULL, int, long, new, private, protected, public, short, 
                signed, static, volatile, String, void, true, unsigned, word, array, 
                sizeof, dynamic_cast, typedef, const_cast, struct, static_cast, union, 
                friend, extern, class, reinterpret_cast, register, explicit, inline, 
                _Bool, complex, _Complex, _Imaginary, atomic_bool, atomic_char, 
                atomic_schar, atomic_uchar, atomic_short, atomic_ushort, atomic_int, 
                atomic_uint, atomic_long, atomic_ulong, atomic_llong, atomic_ullong, 
                virtual, PROGMEM,
                },  
% 
%
  %%% Keyword Color Group 3 %%%  (called KEYWORD1 by arduino)
  keywordstyle=[3]\bfseries\color{arduinoOrange},
  keywords=[3]{  % add built-in functions as a 3rd group
                Serial, Serial1, Serial2, Serial3, SerialUSB, Keyboard, Mouse,
                },      
%
%
  %%% Keyword Color Group 4 %%%  (called KEYWORD2 by arduino)
  keywordstyle=[4]\bfseries\color{arduinoOrange},
  keywords=[4]{  % add more built-in functions as a 4th group
                abs, acos, asin, atan, atan2, ceil, constrain, cos, degrees, exp, 
                floor, log, map, max, min, radians, random, randomSeed, round, sin, 
                sq, sqrt, tan, pow, bitRead, bitWrite, bitSet, bitClear, bit, 
                highByte, lowByte, analogReference, analogRead, 
                analogReadResolution, analogWrite, analogWriteResolution, 
                attachInterrupt, detachInterrupt, digitalPinToInterrupt, delay, 
                delayMicroseconds, digitalWrite, digitalRead, interrupts, millis, 
                micros, noInterrupts, noTone, pinMode, pulseIn, pulseInLong, shiftIn, 
                shiftOut, tone, yield, Stream, begin, end, peek, read, print, 
                println, available, availableForWrite, flush, setTimeout, find, 
                findUntil, parseInt, parseFloat, readBytes, readBytesUntil, readString, 
                readStringUntil, trim, toUpperCase, toLowerCase, charAt, compareTo, 
                concat, endsWith, startsWith, equals, equalsIgnoreCase, getBytes, 
                indexOf, lastIndexOf, length, replace, setCharAt, substring, 
                toCharArray, toInt, press, release, releaseAll, accept, click, move, 
                isPressed, isAlphaNumeric, isAlpha, isAscii, isWhitespace, isControl, 
                isDigit, isGraph, isLowerCase, isPrintable, isPunct, isSpace, 
                isUpperCase, isHexadecimalDigit, 
                },      
%
%
  %%% Set Other Colors %%%
  stringstyle=\bfseries\color{arduinoDarkBlue},    
  commentstyle=\color{arduinoGrey},    
%          
%   
  %%%% Line Numbering %%%%
  numbers=none,%left,                    
  numbersep=5pt,                   
  numberstyle=\color{arduinoGrey},    
  %stepnumber=2,                      % show every 2 line numbers
%
%
  %%%% Code Box Style %%%%
  breaklines=true,                    % wordwrapping
  tabsize=2,
  basicstyle=\ttfamily  
}

%%%%%%%%%%%%%%%%%%%%%%%%%%%%%%%%%%%%%%%%%%%%%%%%%%%%%%%%%%%%%%
%%%%%%%%%%%%%%%%%%%%%%%%%%%%%%%%%%%%%%%%%%%%%%%%%%%%%%%%%%%%%%
%algorithmes en langage courant
%%%%\usepackage[linesnumbered,french]{algorithm2e}
%exemple :
%\normalem %%%% disable auto underline
%\DontPrintSemicolon %n'affiche pas les ; en fin de ligne
%%\SetAlgoNoLine %Enlève des filets verticaux et met les blocs de "fin"
%\SetAlgoVlined %Conserve les filets verticaux avec crochet retour et enlève les blocs de "fin" 
%%\SetAlgoLined  %Conserve les filets verticaux et met les blocs de "fin"
%%\RestyleAlgo{ruled} %boxed, boxruled, ruled and algoruled.
%
%\hspace*{2em}\textbf{Algorithme MaxCompatible(S)}
%
%\hspace*{3em}\begin{algorithm}[H]
%$S' \leftarrow \varnothing$\;
%Trier les activités de S par durée croissante\;
%\Pour{$i$ de 1 à $\abs{S}$}{
%\Si{l'activité $i$ est compatible avec les activités de $S'$}{
%Ajouter l'activité $i$ à $S'$
%}
%}
%Afficher $S'$
%\end{algorithm}
%\ULforem %%%% enable auto underline
%%%%%%%%%%%%%%%%%%%%%%%%%%%%%%%%%%%%%%%%%%%%%%%%%%%%%%%%%%%%%%
%%%%%%%%%%%%%%%%%%%%%%%%%%%%%%%%%%%%%%%%%%%%%%%%%%%%%%%%%%%%%%
%\usepackage{cclicenses}
%exemple : Licence : \cc - \ccby - \ccnc}

\usepackage[scale=1.3]{ccicons}
%Licence : \ccbyncsaeu


%%%%%%%%%%%%%%%%%%%%%%%%%%%%%%%%%%%%%%%%%%%%%%%%%%%%%%%%%%%%%%
%%%%%%%%%%%%%%%%%%%%%%%%%%%%%%%%%%%%%%%%%%%%%%%%%%%%%%%%%%%%%%

\usepackage{fontawesome}

%\faYoutubePlay
%\newcommand{\videoyoutube}[1]{
%\raisebox{-3pt}{\includegraphics[scale=0.1]{youtube}}\,\url{#1}
%}
\newcommand{\videoyoutube}[1]{\faYoutubePlay\ \url{#1}}


\newcommand{\voir}[1]{\textcolor{purple}{\faBook\ voir #1}}

\newcommand{\exoschoisis}[1]{\faBook\ Exercices associés~:\ #1}

\newcommand{\augustin}[1]{\includegraphics[width=#1\linewidth]{Logo_augustin_thierry}}
\newcommand{\acad}[1]{\includegraphics[width=#1\linewidth]{logo_acad}}

\newcommand{\geogebra}[1]{\href{#1}{\raisebox{-0.35em}{\includegraphics[scale=0.3]{geogebra}}\ #1\ \raisebox{-0.35em}{\includegraphics[scale=0.3]{geogebra}}}}

\newcommand{\activite}[2][]{\textcolor{purple}{\faBook\ Activité#1 d'introduction : Situation#1 #2}}
